% Conclusion et perspectives



Tout corps de rapport doit se terminer par une section ou un chapitre "conclusions 
et perspectives". 

Il est courant que le lecteur commence même par lire ce chapitre avant tout le reste pour directement évaluer l'intérêt de regarder plus ou pas, c'est donc une partie à particulièrement soigner.

%********************
\section{Conclusions}
%********************

Dans cette rubrique on doit trouver les conclusions scientifiques ou techniques, qui sont étayées par les résultats obtenus. 

Il faut rappeler les objectifs du projet, et en quelques lignes les méthodologies employées. 

Il faut ensuite rappeler les principaux résultats et les conclusions qu'on en tire. Il faut parler de fait, et éviter les supputations.

Enfin il faut insister sur les aspects originaux du travail et son apport personnel\footnote{On utilisera le passif ou le "on". Il faut éviter le "je" car il est très rare que vous ayez fait le travail tout seul, sans aucune aide extérieure.}

%*********************
\section{Perspectives}
%*********************
 
Il faut démarrer sur des perspectives directes, à courte échéance, qui correspondent à une suite directe du travail accompli.

Ensuite on peut laisser plus de place à l'imagination en donnant des perspectives plus lointaines, mais réalistes. Cela peut devenir de nouvelles orientations du projet en cours.

Les perspectives sont souvent les parties les plus lues en détail car c'est l'aboutissement du travail et une potentielles sources d'inspiration pour continuer le projet.

% *********************************
\section{Conclusions personnelles}
% **********************************

Cette section n'est à mettre que pour un rapport de stage ou de projet
et uniquement à destination des universitaires.

Dans les rapports d'entreprise on n'a pas ce genre de section, sauf si cela
est demandé explicitement.

Il faut rester assez neutre, en rappelant les côtés positifs du stage, tout en 
énonçant les difficultés rencontrées, et en indiquant finalement 
ce qui a été appris, et les nouvelles compétences acquises.



%***********************************
\section{**Conclusion, perspectives}
% **********************************

Quelque soit la rédaction, un document contient nécessairement une introduction et une conclusion.

Pour un compte-rendu, cela est l'occasion de rappeler l'intérêt du TP, ce qu'on a pu apprendre, les difficultés rencontrées et donner par exemple son avis personnel.

On peut ajouter des perspectives.

Pour un compt-rendu,   10 à 15 lignes suffisent amplement.


\vspace{1cm} 
 
La qualité du rapport démontre votre investissement. Un étudiant qui n'a rien compris ou rien fait en TP ne pourra jamais faire un bon rapport. 
Un beau et bon rapport mettra en valeur vos compétences, vos connaissances 
et votre capacité de travail et de compréhension. 
 
\Important{Enfin tout rapport doit contenir une bibliographie, avec des références contenant toutes les informations pour retrouver le document, comme ce qui est fait dans l'exemple ci-dessous.
Cela peut être juste le cours du professeur ou un livre.}
